\chapter{HydrologicalTwinAlphaSeries} \label{chap-proto-hydrotwin}

\phantomsection
\fancyhead[LE]{\thepage} 
\fancyhead[RE]{\textbf{HydrologicalTwinAlphaSeries}}
\fancyhead[RO]{\thepage} 
\fancyhead[LO]{\textbf{HydrologicalTwinAlphaSeries}}
\renewcommand{\headrulewidth}{1pt}

\setlength{\parindent}{0cm}
\setlength{\parskip}{9pt}

% --------------------------------------------------------------

This chapter documents the public \texttt{HydrologicalTwinAlphaSeries} repository, currently published as the standalone Python package \texttt{hydrological\_twin\_alpha\_series}. The alpha series is the public, read-only implementation of the \hydrot\ concepts for \cwv\ and related frontends. It preserves the monolithic backend façade and the six-layer conceptual decomposition of \hydrot\ while restricting the exposed operations to visualization, extraction, comparison, and persistence workflows compatible with open public deployment.

\section{Scope and Positioning}

\subsection{Role within the \hydrot\ programme}

\texttt{HydrologicalTwinAlphaSeries} is not a separate ontology; it is the public implementation branch of the \hydrot\ programme. It gives operational form to the same conceptual commitments already introduced in Part~I:

\begin{itemize}
    \item a single backend façade exposed to client applications,
    \item explicit compartment-, layer-, and observation-level accessors,
    \item deterministic and serializable backend state,
    \item strict separation between the public read-only backend and the future authenticated, calibration-capable \hydrot\ implementation.
\end{itemize}

The priority is therefore given to the name \hydrot\ for the global concept, while \texttt{HydrologicalTwinAlphaSeries} designates the current public implementation series.

\subsection{Relationship to the full \hydrot\ architecture}

The alpha series keeps the architectural grammar of \hydrot\ but implements only the subset compatible with public distribution:

\begin{itemize}
    \item \textbf{Model layer}: compartment, mesh, layer, and configuration objects are exposed through serializable accessors.
    \item \textbf{Data layer}: binary simulation outputs and observations can be read through the backend façade.
    \item \textbf{Estimation layer}: comparison and performance statistics are available for public evaluation workflows.
    \item \textbf{Analysis layer}: extraction and temporal/spatial post-processing remain backend responsibilities.
    \item \textbf{Cartographic layer}: rendering-oriented data products remain consumable by \cwv.
    \item \textbf{Registry / persistence}: the public series currently relies on deterministic package state and pickle persistence rather than the full authenticated registry model.
\end{itemize}

In contrast, the full \hydrot\ target still includes private-key authentication, parameter editing, simulation execution, self-fitting, frequency-domain calibration, and Bayesian inference workflows. The alpha package should thus be read as the public operational backend of \hydrot\, not as a replacement for the complete framework.

% --------------------------------------------------------------

\section{Standalone Python Package Architecture}

\subsection{Repository and package identity}

The public repository is distributed independently from the QGIS plugin under the package name \texttt{hydrological\_twin\_alpha\_series}. The package is installed as a normal Python dependency and then consumed by \cwv\ through an editable install from a checked-out repository or submodule.

This is an important evolution with respect to the earlier embedded prototype: the backend is no longer documented as an internal \texttt{cawaqsviz.ht} subpackage. It is now a standalone backend package whose documented public import surface is the canonical package entrypoint:

\begin{verbatim}
from HydrologicalTwinAlphaSeries import (
    ConfigGeometry,
    ConfigProject,
    HydrologicalTwin,
)
\end{verbatim}

\subsection{Current package layout}

The current implementation is organized around a standalone package root that keeps the historical CaWaQS domain modules together with a dedicated \texttt{ht/} façade package:

\begin{verbatim}
src/hydrological_twin_alpha_series/
    __init__.py
    Compartment.py
    Extraction.py
    Manage.py
    Mesh.py
    Observations.py
    Renderer.py
    Vec_Operator.py
    config/
        __init__.py
        ...
    ht/
        __init__.py
        hydrological_twin.py
        api_types.py
        persistence.py
    tools/
        ...
\end{verbatim}

This organization is conceptually consistent with \hydrot\:
\begin{itemize}
    \item the package root retains the domain objects and legacy scientific operators;
    \item the internal \texttt{ht} subpackage isolates the monolithic \texttt{HydrologicalTwin} façade and its private delegation paths;
    \item configuration and utility concerns remain separated from the façade itself.
\end{itemize}

The repository still contains internal modules and historical scientific operators, but the White Book now treats only the canonical monolithic class and its seven macro methods as supported public API. Transitional wrappers and deleted helper shims are not part of the intended user contract.

\subsection{Backend façade}

The central public class is now \texttt{HydrologicalTwin}, documented from the canonical package entrypoint:

\begin{verbatim}
from HydrologicalTwinAlphaSeries import HydrologicalTwin
\end{verbatim}

The class remains the monolithic backend façade for \cwv. Internally it still delegates to the historical CaWaQS operators, but the public contract is now expressed through the canonical seven-method macro API:
\texttt{configure}, \texttt{load}, \texttt{describe}, \texttt{extract}, \texttt{transform}, \texttt{render}, and \texttt{export}. This preserves the six-layer ontology of \hydrot\ while removing the earlier transitional public helper surface.

% --------------------------------------------------------------

\section{Public Python API}

\subsection{Construction}

The minimal public construction pattern currently documented by the repository is the following:

\begin{verbatim}
from HydrologicalTwinAlphaSeries import (
    ConfigGeometry,
    ConfigProject,
    HydrologicalTwin,
)

config_geom = ConfigGeometry.fromDict({...})
config_proj = ConfigProject.fromDict({...})

twin = HydrologicalTwin(
    config_geom=config_geom,
    config_proj=config_proj,
    out_caw_directory=config_proj.cawOutDirectory,
    obs_directory=config_proj.obsDirectory,
)
\end{verbatim}

The constructor stores the geometry and project configuration, output and observation roots, an optional temporary directory, and optional metadata. It also instantiates the service objects reused from the CaWaQS scientific stack:\ \texttt{Manage.Temporal()}, \texttt{Manage.Spatial()}, and \texttt{Manage.Budget()}.

\subsection{Canonical seven-method macro API}

The documented public surface is now intentionally narrow. User-facing workflows are expressed through seven macro entrypoints:

\begin{itemize}
    \item \texttt{configure}: register project configuration, runtime options, and access mode;
    \item \texttt{load}: materialize the twin state from declared inputs and cached artefacts;
    \item \texttt{describe}: inspect available compartments, layers, variables, observations, and loaded resources;
    \item \texttt{extract}: retrieve normalized scientific products such as time series, spatial maps, observations, and aquifer outcropping subsets;
    \item \texttt{transform}: apply temporal, spatial, budget, or statistical operators to extracted products;
    \item \texttt{render}: create public-facing visualizations, including simulation-observation plots;
    \item \texttt{export}: persist canonical outputs such as CSV tables, figures, reports, or cache artefacts.
\end{itemize}

These macro methods are the only operations that should be used in user-guide examples. Legacy transition helpers such as \texttt{register\_compartment}, \texttt{describe\_api\_facade}, \texttt{build\_*\_gdf}, and direct \texttt{render\_sim\_obs\_*} helpers are no longer documented as public API.

A representative public workflow is therefore:

\begin{verbatim}
twin.configure(config_geom=config_geom, config_proj=config_proj)
twin.load(kind="public_backend")
summary = twin.describe(kind="compartments")
series = twin.extract(kind="timeseries", compartment="HYD", variable="discharge")
regime = twin.transform(kind="temporal_aggregate", source=series, operator="monthly_mean")
figure = twin.render(kind="sim_obs_pdf", source=regime)
twin.export(kind="csv", source=regime, destination="/path/to/regime.csv")
\end{verbatim}

This documentation change is important: the current public alpha implementation is no longer presented as a collection of transitional helper calls. It is centred on a single \texttt{HydrologicalTwin} façade whose canonical macro API stabilizes the backend contract for \cwv, scripts, and future services.

\subsection{Migration from deleted helper methods}

The final cleanup of the alpha repository removed several public transition helpers while keeping their behaviour available behind the canonical macro API. The corresponding migration guidance is:

\begin{center}
\begin{tabular}{p{0.33\textwidth} p{0.55\textwidth}}
\textbf{Legacy helper usage} & \textbf{Canonical replacement} \\
\hline
Spatial maps and \texttt{build\_*\_gdf} wrappers & \texttt{extract(kind="spatial\_map", ...)} \\
Aquifer outcropping helper & \texttt{extract(kind="aquifer\_outcropping", ...)} \\
Simulation-observation PDF rendering & \texttt{render(kind="sim\_obs\_pdf", ...)} \\
Simulation-observation interactive rendering & \texttt{render(kind="sim\_obs\_interactive", ...)} \\
Ad hoc API-surface description helpers & \texttt{describe(kind=..., ...)} \\
\end{tabular}
\end{center}

This migration rule is intentionally simple: all remaining internal delegation paths stay inside the class, while the documented user surface remains the seven-method canonical API.

% --------------------------------------------------------------

\section{Persistence and Public Reproducibility}

\subsection{Lightweight public provenance and pickle-based persistence}\label{p2c2_ssGitRec}

The standalone package currently provides a lightweight public reproducibility mechanism aligned with the canonical macro workflow:

\begin{verbatim}
twin.export(kind="pickle", destination="/path/to/cache/twin_snapshot.pkl")
twin.load(kind="pickle", source="/path/to/cache/twin_snapshot.pkl")
\end{verbatim}

The persisted payload remains intentionally lighter than the target \texttt{GitSynchroRecord}-centred registry of the full \hydrot\ architecture, but it still supports reproducible public workflows, bancarisation, and fast session recovery.

\subsection{Interpretation within the White Book}

From the standpoint of the White Book, this means that the public alpha series already validates three important principles:

\begin{enumerate}
    \item the backend state can be externalized and restored deterministically;
    \item the façade can remain monolithic without collapsing conceptual layering;
    \item public reproducibility can be achieved before private authenticated capabilities are added.
\end{enumerate}

% --------------------------------------------------------------

\section{Integration with \cwv}

\subsection{Current integration mode}

The public repository is designed to be integrated into \cwv\ as an external checkout, typically through a Git submodule mounted at:

\begin{verbatim}
external/HydrologicalTwinAlphaSeries/
\end{verbatim}

The plugin environment then installs the backend in editable mode. This supersedes the earlier documentation pattern in which the backend was presented as source files copied directly inside the \texttt{cawaqsviz} package.

\subsection{Architectural consequence}

This standalone packaging strengthens the architectural separation already defended in this White Book:

\begin{itemize}
    \item \cwv\ remains the visualization client,
    \item \texttt{HydrologicalTwinAlphaSeries} remains the scientific backend,
    \item the long-term \hydrot\ programme remains the conceptual and authenticated superset.
\end{itemize}

In other words, the public alpha repository now makes explicit in code what the White Book has argued in theory: the backend should be a reusable hydrological twin service, not an implementation detail hidden inside the GUI repository.

% --------------------------------------------------------------

\section{Conclusion}

The publication of \texttt{HydrologicalTwinAlphaSeries} clarifies the implementation trajectory of \hydrot. The public alpha series is now a standalone Python backend package exposing a canonical \texttt{HydrologicalTwin} façade organized around the seven macro methods \texttt{configure}, \texttt{load}, \texttt{describe}, \texttt{extract}, \texttt{transform}, \texttt{render}, and \texttt{export}, while preserving conceptual continuity with the six-layer \hydrot\ architecture.

This update therefore replaces the earlier description of an embedded \texttt{HydroTwin.py} prototype by a more accurate picture: \texttt{HydrologicalTwinAlphaSeries} is the current public implementation branch of \hydrot\ for \cwv\ and related frontends, whereas the complete \hydrot\ target remains the authenticated, calibration-capable digital twin framework described in the rest of this White Book.
